\section{What the Sequent Calculus Rules Out: Speculation}
\label{sec:speculation}

Eisner and Blatz~\cite{eisner-blatz-2007-parsing} introduce \emph{speculation}
as a program transformation for weighted logic programs: if a subcomputation
is independent of some loop variable, hoist it out.  In the SC literature this
appears as ``loop-invariant code motion for logic programs.''

The sequent-calculus framework gives us the vocabulary to characterise exactly
why speculation cannot be expressed as a generalisation step.

\subsection{The reduction}

In our setting, speculation would correspond to: given configs $j_1, j_2$,
find a subterm $s$ of $j_1$ independent of dropped variables $V = \fv(j_2)
\setminus \fv(j_1)$, project the typing context to only $\fv(s)$, and use
the projected config as the generalised node.

The implementation (\texttt{SpeculationGen.v}) detects independent subterms
via a syntactic free-variable check, projects the context using the existing
\texttt{canon\_jTy}, and produces a $\mathsf{gen\_result}$.

\subsection{The negative result}

\begin{theorem}[Speculation redundancy]
For canonical configurations $j_1, j_2$, if $\mathsf{generalize\_speculation}\ j_1\ j_2$
produces a $\mathsf{gen\_result}$ and the resulting fold is CIU-sound, then
$j_1$ and $j_2$ have the same term --- making speculation redundant with
anti-unification.
\end{theorem}

\begin{proof}[Proof sketch]
Let the gen\_result have term $s$ (the independent subterm) and sigma
$\mathsf{subterm\_sigma}\ s = \mathsf{map}\ \mathtt{tVar}\ (\fv(s))$.
The fold replaces a vertex with
$\mathsf{apps}(\mathsf{residual}(\mathsf{gen\_j}), \mathsf{sigma})$.
By \texttt{canon\_subterm\_roundtrip}, this reconstructs $s$, not the
original term $t$ (which contains $s$ plus additional structure).
For the fold to recover $t$, we must have $t = s$, i.e.\ the whole term
is independent.  But then the two configs have identical terms (after
canonicalisation), and anti-unification handles them.
\end{proof}

\noindent
The consequence is a \emph{characterisation}: in canonical configurations,
speculation as a generalisation step adds no power beyond anti-unification.
The same conclusion follows from the proof theory: anti-unification
abstracts over differing subterms (identity expansion), while projection
dropping variables requires a splitting step (cut with a fresh variable)
that is a different proof-theoretic operation.

\subsection{What this means}

This negative result is a product of the sequent-calculus methodology.
Without the precise vocabulary of context projection and cut formulas,
the distinction between generalisation (identity, extending context)
and splitting (cut, restricting context) would remain blurred.
The framework does not merely enable new techniques --- it also
identifies which proposed techniques cannot work, saving the effort
of pursuing dead ends.
